A resistência de um resistor, $R$, é feita passando várias correntes, $I$, através dele e medindo as quedas de tensão correspondentes, $V$, e correntes com medidores analógicos imprecisos. Os dados resultantes são: $x_i = V$ (V): 10, 20, 30, 40, 50, 60, 70, 80; $y_i = I$ (A): 0.8, 1.1, 2.5, 4.2, 4.3, 4.7, 5.8, 6.4. (a) Que tipos de erro serão encontrados nos dados? (b) Supondo que $V = IR$, plote os dados (manualmente!) e estime visualmente (globo ocular) a linha de melhor ajuste para esses dados.

